% Generated draft: requires mathematical and source review.
\documentclass[11pt]{article}
\usepackage[margin=25mm]{geometry}
\usepackage{fontspec}
\setmainfont{Latin Modern Roman}
\usepackage{amsmath,amssymb,mathtools,unicode-math}
\setmathfont{Latin Modern Math}
\usepackage{enumitem,needspace,xurl,hyperref}
\setlength{\parindent}{0pt}
\setlength{\parskip}{5pt}
\setlength{\emergencystretch}{3em}
\newcommand{\N}{\mathbb{N}}
\newcommand{\Z}{\mathbb{Z}}
\newcommand{\Q}{\mathbb{Q}}
\newcommand{\R}{\mathbb{R}}
\newcommand{\C}{\mathbb{C}}
\newcommand{\F}{\mathbb{F}}
\newcommand{\A}{\mathbb{A}}
\newcommand{\PP}{\mathbb{P}}
\newcommand{\E}{\mathbb{E}}
\newcommand{\eps}{\varepsilon}
\newcommand{\dd}{\,\mathrm{d}}
\begin{document}
\section*{UnsolvedMath ID: 6200019}
\textbf{Source problem code:} catalogue entry
\subsection{Definitions and mathematical statement}
A CAT(0) cube complex is a cube complex with its standard piecewise Euclidean cubical length metric satisfying the CAT(0) comparison inequality. A geometric group action is an action by isometries that is proper and cocompact. Use the usual proper geodesic-space setting for its visual boundary.

For a proper CAT(0) space $X$, the visual boundary $\partial_\infty X$ consists of asymptotic classes of geodesic rays, where two rays are equivalent if they remain at bounded distance. It has the cone topology, described from a basepoint by uniform convergence of rays on bounded parameter intervals.

Let a group $G$ admit geometric actions on two CAT(0) cube complexes $X_1$ and $X_2$. Must
\[
\partial_\infty X_1\cong\partial_\infty X_2
\]
as topological spaces? Equivalently, is the homeomorphism type of the visual boundary determined by the group within this cubical class of geometric models? The source's equality notation is interpreted as homeomorphism type, not literal equality of boundary sets. It does not explicitly require the homeomorphism to be $G$-equivariant. Nonuniqueness examples for more general CAT(0) spaces, mentioned in the summary, are not by themselves counterexamples within the stated cubical class.
\subsection{Short English statement}
If the same group acts properly and cocompactly on two CAT(0) cube complexes, must their visual boundaries be homeomorphic?
\subsection{Sources}
\begin{itemize}
\item Misha Kapovich, \emph{Problems on Boundaries of Groups and Kleinian Groups} (2005), Problem 19, PDF page 7, attributed to Dani Wise, \url{https://www.math.ucdavis.edu/~kapovich/EPR/problems.pdf}; accessed 2026-07-29, PDF-text-normalized extraction.
\item Local registry \texttt{data/unsolvedmath/problems.json}, permanent UnsolvedMath ID $6200019$, source code catalogue entry, \emph{Boundaries of Groups and Kleinian Groups --- Problem 19}. Dataset attribution: Ulam AI and the UnsolvedMath Contributors.
\end{itemize}
\end{document}
